% GDZ PPN236018647,native536;HKUST537.
\clearpage\SourceLeaf{536}\ArithmeticHead{507}{R}{1796 OCTOB. 1.—1796 OCT. 21.}
\NumberedItem{40}
\indent Aequationis
\[
x^p-1=0
\]
\noindent radices per integros multiplicatae aggregatae cifram producere non possunt.
\par\noindent\hfill[1796] \(\odot\) Oct. 9. Brunsv[igae]

\medskip\begingroup\small
\indent Wenn \(p\) eine Primzahl ist, so bedeutet die Aussage, daß die primitiven Wurzeln von \(x^p-1=0\) mit ganzen Zahlen multipliziert nicht die Summe Null geben können, die Irreduzibilität der Kreisteilungsgleichung; siehe den art. 341. der \textit{Disquisitiones arithmeticae}, Werke I, S. 417. Für einen zusammengesetzten Exponenten vergleiche man die Tagebuchaufzeichnung Nr. 136 vom 12. Juni 1808.
\par\noindent\hfill\textsc{Bachmann.}
\par\endgroup

\NumberedItem{41}
\indent Quaedam sese obtulerunt de multiplicatoribus aequationum, ut certi termini eliciantur, quae praeclara pollicentur.
\par\noindent\hfill[1796] \(\odot\) Oct. 16. Bruns[vigae]

\medskip\begingroup\small
\indent Unter multiplicatores versteht \textsc{Gauss} hier wohl Transformationen, möglicherweise \textsc{Tschirnhausen}sche, also solche von der in der Nr. 34 betrachteten Form. Man hätte dann die vorliegende Aufzeichnung auf die Beseitigung gewisser Gleichungskoeffizienten durch \textsc{Tschirnhausen}sche Transformationen zu beziehen. Die Wendung »praeclara pollicentur« scheint anzudeuten, daß die Frage nach der Auflösbarkeit der allgemeinen Gleichung durch Wurzelgrößen das Ziel dieser Untersuchungen gebildet habe, und die beiden folgenden, für sich nicht zu deutenden Nummern 42, 43 können Fortschritte auf dem hier eingeschlagenen Wege ankündigen. Die rote Unterstreichung in der Handschrift würde dann auf die in der Bemerkung zu der Nr. 37 genannten Veröffentlichungen in der \textit{Dissertation} und in den \textit{Disquisitiones arithmeticae} zu beziehen sein.
\par\noindent\hfill\textsc{Loewy.}
\par\endgroup

\NumberedItem{42}
\indent Lex detecta: quando et demon[stra]ta erit systema ad perfectionem evexerimus.
\par\noindent\hfill[1796] Oct. 18. Brunsv[igae]

\NumberedItem{43}
\indent Vicimus GEGAN.
\par\noindent\hfill[1796] Oct. 21, Bruns[vigae]
\par\vfill\noindent\hfill64\textsuperscript{*}
